\section{Four original quadratic families on the same actual host}
\label{sec:cubic-quadratic-family-transfer}

The auxiliary seed and native support of AC-G123 depend on a residual
tuple over their actual host.  They do not require an original Type II
endpoint in each weight.  We prove that the same support applies to four
original quadratic families.  Two of them have a ramified quadratic
phase at the coefficient prime.  Their selected local components are
checked separately; equality of their residual tuples is not used to
identify their integral local deformation rings.

\paragraph{Characters, original models, and retained inputs.}
Put
\[
 p=269,\qquad K=\mathbb Q(\sqrt{-3}),\qquad
 \psi_0=\chi_{-4},\qquad \delta=\chi_5,\qquad \tau=\chi_{269}.
 \tag{G124.1}
\]
Here the notation denotes the primitive quadratic characters of the
indicated fundamental discriminants; \(\tau\) is even.  The actual
cyclotomic character is \(c_p\), with
\(c_p(\operatorname{Fr}^{\mathrm{geom}}_w)=q_w^{-1}\).
For \(s,t\in\{0,1\}\) define
\[
 \begin{aligned}
 \psi_{s,t}&=\psi_0\delta^s\tau^t,\qquad
 N_{s,t}=4\cdot5^s\cdot269^t,\\
 e&\equiv267-134t\pmod {268},\quad e>0,\qquad
 \gamma_{s,t,e}=\psi_{s,t}c_p^e.
 \end{aligned}
 \tag{G124.2}
\]
The four primitive data are
\[
\begin{array}{c|c|r|c|c}
 (s,t)&\psi_{s,t}&N_{s,t}&e\pmod {268}&\operatorname{comp}_p\psi_{s,t}\cr
 (0,0)&\chi_{-4}&4&267&1\cr
 (1,0)&\chi_{-20}&20&267&1\cr
 (0,1)&\chi_{-1076}&1076&133&\tau\cr
 (1,1)&\chi_{-5380}&5380&133&\tau
\end{array}
\tag{G124.3}
\]

Fix the original coefficient DVR \(\mathcal O_0\), its fraction field
\(E_0\), residue field \(k_0\), and a label \(j\in\{1,2\}\).
For each datum in (G124.2), take its own actual normalized cubic ring
\(S_{j,s,t,e}\), selected local prime \(P_{*,s,t,e}\), and original
model
\[
 (A_0,\mathcal A,x,q_0:A_0\longrightarrow S_{j,s,t,e},\mathcal Q)
 \tag{G124.4}
\]
as in AC-G65--AC-G66.  A repeated symbol in (G124.4) denotes the
corresponding object for that datum, not an identification of the
different original rings.  Let \(\mathscr B_{0,j,s,t,e}\) be its
complete degree-zero subring and \(L_{j,s,t,e}\) its actual
weight-minus-one cotangent length.

Retain the exact original coefficient, \(\mathrm{GD}/\mathrm{Dat}\),
marked-element and selected-component interfaces of AC-G65--AC-G68,
and their specified published local inputs.  Retain
\(\mathrm{CohInp}\), \(\mathrm{LemmaP}\),
\(\mathrm{LemmaM}\), \(\mathrm{LemmaM'}\), and the old
\(\mathrm{AutomorphyUnramifiedCoefficientsAlt}\) supply.  In
particular retain
\[
 e_{\mathcal O_0}
       =\operatorname{length}_{\mathcal O_0}(\mathcal O_0/p)=1
 \tag{G124.5}
\]
where that original modularity supply requires it.
Keep precisely the native, attachment, ordinary-control, transfer,
cyclic-base-change, class-field, duality, adjoint and counting inputs
listed in AC-G123.  Thus \(\mathrm{ModU}\), \(\mathrm{LamEx}\),
\(\mathrm{LevEx}\) and \(\mathrm{ProModularRingExists}\) remain
explicit native supplies, not consequences of A alone.
The theorem below supplies the actual host support from AC-G123; it
does not add a separate support or original endpoint-existence premise.

\begin{theorem}[AC-G124-family-transfer: four actual model families]
\label{AC-G124-family-transfer}
Under these retained inputs, for fixed \(\mathcal O_0\) and \(j\),
one may choose the final host \(H/F\), host coefficients, residual
tuple rings \(B\subset A\subset R\), native modular ideal \(J_3\),
map \(\Phi_3\), joint-kernel weight ideal \(\lambda_\theta\), and
support ideals of AC-G123 once for all four rows of (G124.3) and all
their indicated positive weights.  With
\[
 I=\Phi_3(J_3)B,\qquad I_{\mathrm{off}}=I\cap M_E,
 \tag{G124.6}
\]
these actual ideals satisfy
\[
 \begin{aligned}
 I&\subset\mathfrak q\cap B
       &&\text{for every Type I prime of }A,\\
 I_{\mathrm{off}}&\subset\mathfrak q\cap B
       &&\text{for every off-triangular prime of }A,\\
 \sqrt{I_{\mathrm{off}}+\lambda_\theta}&=\mathfrak m_B.
 \end{aligned}
 \tag{G124.7}
\]
Every original model (G124.4) has \(\mathscr B_{0,j,s,t,e}\) finite
over \(\mathcal O_0\) and satisfies the five literal conditions
\(\mathrm{M1}\)--\(\mathrm{M5}\).  Every two-sided
characteristic-zero local point is Type I and rigid.

For a fixed \((s,t,e)\), put
\[
 \begin{split}
 \alpha_j&=\chi^{-j}\psi_{s,t}^{-1}c_p^{-e}|_{G_K},\\
 \mathcal S_j&=\ker\Bigl(
 H^1(G_K,(E_0/\mathcal O_0)(\alpha_j))\\
 &\qquad\longrightarrow
 \prod_{v\mid p}H^1(G_{K_v},(E_0/\mathcal O_0)(\alpha_j))\\
 &\qquad\qquad\times\prod_{v\nmid p}H^1(I_v,(E_0/\mathcal O_0)(\alpha_j))
 \Bigr),\\
 s_j&=\operatorname{length}_{\mathcal O_0}(\mathcal S_j)/e_{\mathcal O_0},
 \qquad Y=N_{s,t}(e+1),\qquad b=121^e+1+121^{-e}.
 \end{split}
 \tag{G124.8}
\]
Their lengths are finite.  At geometric Frobenius \(\sigma\) above
the inert rational prime \(11\), a finite family \(X_j\) of actual
Type I representations satisfies
\[
 s_j\le L_{j,s,t,e}/e_{\mathcal O_0}
 \le 2\sum_{\rho\in X_j}
           v_{269}(\operatorname{tr}\rho(\sigma)-b),
 \qquad \#X_j\le C_1Y^{A_1}.
 \tag{G124.9}
\]
The constants are those of AC-G112-count.  For any nonempty set
\(J_0\subset\{1,2\}\) with these original inputs, one nonzero
integer \(\mathcal N_{s,t,e,J_0}\) satisfies
\begin{align}
 \sum_{j\in J_0}L_{j,s,t,e}/e_{\mathcal O_0}
       &\le2v_{269}(\mathcal N_{s,t,e,J_0}),\notag\\
 \log|\mathcal N_{s,t,e,J_0}|
       &\le5C_1\#J_0\,Y^{A_1+1}\log11,\notag\\
 \left(\sum_{j\in J_0}L_{j,s,t,e}/e_{\mathcal O_0}\right)\log269
       &\le10C_1\#J_0\,Y^{A_1+1}\log11.
 \tag{G124.10}
\end{align}
These bounds also hold with \(s_j\) in place of the corresponding
normalized cotangent length.  The two labels may use separate hosts;
the common integer uses their common original-field classifier.
There is no change of \(\mathcal O_0\) or its valuation.
\end{theorem}

\begin{proof}
\proofstep{1. The original arithmetic data.}
The four numbers \(-4,-20,-1076,-5380\) are fundamental
discriminants.  Indeed each is \(4d\) for a squarefree
\(d\equiv3\pmod4\).  Their primitive characters are exactly the
products in (G124.2), so their conductors are (G124.3).  Each
character is odd of order two.  Every indicated \(e\) is odd and
positive; hence \(\psi_{s,t}(-1)=(-1)^e\).  Also
\[
                2e\equiv266\pmod {268}.
 \tag{G124.11}
\]
Thus the literal \(\mathrm{GD.gen}\), parity, prime-to-\(p\)
character-order and even-modulus conditions hold.
Chinese remaindering in the source definition of
\(\operatorname{comp}_p\) gives exactly \(\tau^t\); its square
is one.  The prime-to-\(p\) part \(\chi P\) in
\(\mathrm{GD.eta}\) is \(\psi_0\delta^s\).
The field cut out by this character and \(\bar c_p\) is unramified
away \(2,5,269\).  It cannot contain \(K\), which is ramified
at \(3\).  This is the literal noncontainment in
\(\mathrm{GD.eta}\), not a condition on the auxiliary modular
level.  Moreover \(3\nmid N_{s,t}\), \(h_K=1\), and
\(p\nmid3h_K\).  The remaining quadratic and imaginary-field
clauses hold for \(K\).  Thus the displayed arithmetic fields of
\(\mathrm{GD}\) have been checked, while each actual marked and
local \(\mathrm{Dat}\) interface remains as specified in the
statement.  All the finite character values here are \(\pm1\);
no new value-field roots are needed for (G124.3).

\proofstep{2. The fixed host and the exact residual equality.}
Apply AC-G123 once to \(\psi_0,e_*=267\), with the fixed
\(\mathcal O_0,j\).  It constructs the support from the actual
AC-G119 auxiliary form, not from an original Type II point.
In its initial host
\[
 H_{\mathrm{ab}}
   =K\mathbb Q(\sqrt{-23})\mathbb Q(\sqrt5),
 \tag{G124.12}
\]
the real field \(E=\mathbb Q(\sqrt5)\) lies in the real part.
It remains in the final \(F=H^+\) after the two actual cyclic
\(269\)-extensions, the selected ray extension, and the final
cyclotomic layer.  Consequently \(\delta|_{G_F}=1\).
Quadratic reciprocity at the prime \(269\), or the unique order-two
character of \(\mathbb F_{269}^{\times}\), gives
\(\bar\tau=\bar c_p^{134}\) on \(G_{\mathbb Q}\).
Therefore, as characters of the full group \(G_F\),
\[
 \bar\gamma_{s,t,e}
    =\bar\psi_0\bar c_p^{e+134t}
    =\bar\psi_0\bar c_p^{267}
    =\bar\gamma_0.
 \tag{G124.13}
\]
This also holds on \(G_H\) and every full local group.
Both inverse endpoint characters consequently agree as well.
In particular the actual host cubic residual tuple is exactly
\[
 (\bar\chi^j\bar\gamma_0,\ 1,\ \bar\chi^j\bar\gamma_0^{-1})
 \tag{G124.14}
\]
for every row, with the same ordering and polarization.

No new prime-to-\(p\) class-group ascent is used here.
The numerical initial four-sector certificate, the two cyclic
\(p\)-rank transports, the selected-place positivity and cup
filling, and the final strict dimension gap are exactly those
already proved for this tuple in AC-G123.
The real degree and the coefficient field degree are not
interchanged.  In particular the comparison constant
\(r_\alpha\le3d/4\) is unchanged.
We do not assert that the enlarged original residual splitting
field for \(\psi_0\delta\) meets \(H\) only in \(K\): it
contains \(E\).  What is used is the actual character equality
(G124.13) after restriction, which supplies the same native datum.

\proofstep{3. Away-kernel and native-field matching.}
At every place away from \(p\), \(\tau\) is unramified.
For a place of \(H\), \(\delta\) is trivial on its entire local
Galois group since \(E\subset H\).  Thus the new original
residual kernel on host inertia is the old joint kernel from
(G124.14).  The modified cubic \(\mathrm{Kill}\) condition of
the actual original model kills this kernel over every test
algebra, before passing to a field or semisimplification.
At the additional rational prime \(5\), this assertion is
particularly explicit: \(\chi\) and \(c_p\) are unramified,
\(\psi_0\) and \(\tau\) are unramified, and
\(\delta|_{G_H}=1\).  Hence all host inertia at \(5\) lies
in the original modified \(\mathrm{Kill}\) set and its image
is one.  No relative splitting assumption at \(5\) is needed.
At the other away places the same calculation gives precisely
the existing host inertia types.

Thus the rational host quotient \(\Gamma_{\bar\gamma_0}\),
the cubic joint-kernel quotient, their full local groups and
all residual representations used by AC-G94, AC-G95, AC-G97
and AC-G106-native are unchanged.  So are their normalized
marked host elements; they need not be identified with the
different original marked elements in (G124.4).
The native \(\kappa_2,\kappa_3\) and their norm/transfer
conditions depend on the host and the cyclotomic character,
not on a choice of lift of (G124.13).  The groups \(\mu\)
are the source's \(p\)-primary torsion preimages in the local
abelianizations.  AC-G123's local fields have no \(\mu_p\),
so the same trivial \(\mathrm{th}\) and exact
\(\mathrm{IsPsi.mu}\) checks apply.  A finite quadratic
phase at \(p\) is not forbidden by this condition.

Consequently the actual native \(\mathrm{Resid}\),
\(\mathrm{LamObj}\), full ordinary functor, \(J_3\), and
\(\Phi_3\) are the existing ones, as are the rational rank-two
corner and its uniform \(M_E\).  The original rank-three and
rank-two auxiliary seeds at \(e_*=267\) already establish
properness, the strong Seed component and the uniform even
support on this same native datum.  There is no new seed
existence or pointwise interpolation step.
AC-G108 therefore gives (G124.7) with these actual ideals.

\proofstep{4. The selected local components at \(p\).}
First, \(126^2\equiv5\pmod {269}\), and the derivative of
\(X^2-5\) at \(126\) is nonzero modulo \(269\).  Hensel's
lemma gives \(\sqrt5\in\mathbb Q_{269}\).  Thus
\(\delta|_{G_{\mathbb Q_{269}}}=1\), not merely on inertia.
For fixed \(e\) in the \(t=0\) progression, the two rows
therefore have identical full local \(\Psi,\tau,\eta,K_v\)
data and the corresponding actual selected local condition.

For \(t=1\), the full local character is different.  We use
its own \(P_{*,s,1,e}\) from AC-G65-local, with the finite
phase \(\tau\) retained.  The hypotheses used in that proof
are the actual local \(\mathrm{Dat}\),
\(\psi^2|_{I_p}=1\), \(e\ge1\), and
\(268\nmid2e\); they all hold here.  Its fixed potentially
crystalline type has inertia
\[
                  (\chi^j\tau,1,\chi^j\tau),
 \tag{G124.15}
\]
and its Hodge--Tate weights are \((-e,0,e)\).
AC-G65-flag supplies, over every commutative local test
algebra with the prescribed residue map, a strict full
\(G_{K_p}\)-stable flag with exact ordered inertia characters
\[
        \chi^j\psi_{s,t}c_p^e,\qquad1,\qquad
                         \chi^j\psi_{s,t}^{-1}c_p^{-e}.
 \tag{G124.16}
\]
It applies also to characteristic-\(269\) DVRs and nonreduced
tests.  This is the selected-component theorem itself; an
arbitrary ordinary component is not substituted for it.

After reduction, (G124.16) is exactly (G124.14) on host
inertia.  On the full finite unit torus the same comparison
also retains the ordinary component: the parallel algebraic
weight changes the endpoint unit factors by
\(\omega(\operatorname{Norm}a)^{\pm(e-267)}\), while
\(\tau^t(a)=\omega(\operatorname{Norm}a)^{134t}\).
The exponents \(134t\pm(e-267)\) are both divisible by
\(268\).  Thus the full finite-torus character is the old
one; no \(p\)-divisible averaging is involved.
The varying integral character and the local Hodge type
still belong to their own original model.

For characteristic-zero points the native arithmetic
characters retain that finite phase.  After the common
\(\kappa_3=c_p^{-1}\) normalization their ordered weights
are \((1-e,1,1+e)\), strictly increasing.  Since here
\(e\ge133\), none of their ratios can equal \(c_p\):
the Hodge--Tate differences have absolute values \(e\) or
\(2e\), and a finite character cannot change them.
The source \(\mathrm{IsWtPart.ph}\) is supplied by the
actual finite phase in (G124.16).  This verifies the
native arithmetic-point interface for these points,
without identifying their integral \(p\)-characters with
the auxiliary seed's characters.

\proofstep{5. Original characteristic-\(p\) curves and M2.}
Take any curve tested by AC-G109 in the actual original
model (G124.4).  Use its same-residue extension
\(V\hookrightarrow k''[[z]]=V'\) and
\(C=k+zV'\), with the host residue field enlarged as
necessary.  The original \(\mathcal O_0\) and its model
are not replaced by host coefficients; all polynomial
identities descend along the injective curve extension.
By Step 3 the restricted representation kills the exact
host away kernel, and by Step 4 it has the required paired
full local flags.  The orthogonal normalization of
AC-G109 therefore gives the actual map \(R\to C\).
In characteristic \(p\), the exact constants in
(G124.16) reduce to (G124.14), so this map kills the
actual joint-kernel ideal \(\lambda_\theta\).

The actual \(I_{\mathrm{off}}\) in (G124.7) now supplies
the curve argument of AC-G109.  Its balanced monomials,
triangular conclusion, class-exponent argument and
finite-index polynomial descent apply unchanged.
At the last descent the original diagonal characters are
\(\chi^j\gamma_{s,t,e},1,\chi^j\gamma_{s,t,e}^{-1}\).
They are not replaced on \(G_K\) by the auxiliary
\(\psi_0c_p^{267}\).  The source polynomial identities
and the generic finite-criterion argument consequently
give M2 for (G124.4), namely the finiteness of
\(\mathscr B_{0,j,s,t,e}\) over \(\mathcal O_0\).
This step makes no map between two different original
deformation rings and no assertion of full Hecke-curve
equality.

\proofstep{6. Characteristic-zero points and M4.}
Every original Type I point restricts to the same host
ordinary functor by Steps 3--4.  Its actual determinant
and all local characters define the native point used
in AC-G110.  The first line of (G124.7), using
\(I=\Phi_3(J_3)B\) itself rather than
\(I_{\mathrm{off}}\), places it on native modular
support.  AC-G110's all-open and full-cyclotomic-kernel
irreducibility proof uses the actual cubic separation
and the nonzero cyclotomic exponents, which hold for
each current \(e\).  The arithmetic-point theorem and
the retained adjoint theorem then give the original
adjoint rigidity by characteristic-zero restriction
and corestriction.  A degree divisible by \(269\)
is invertible in this coefficient field.

For Type II, use AC-G117 with the actual original
quadratic \(\psi_{s,t}\).  That theorem explicitly
allows \(\operatorname{comp}_p\psi=\tau\).
It untwists that phase before using the primitive
conductor-three newform and the all-residue Eisenstein
exclusion.  Thus the original Type II locus is empty
in all four rows; no mixed Type II Bloch--Kato vanishing
is asserted or needed.  The literal tangent-model
argument of AC-G122-M4 gives M4.  AC-G66 supplies M1
and M3 on the original selected models.

\proofstep{7. The test trace, principal level, and original lengths.}
The prime \(11\) is prime to \(3\cdot269N_{s,t}\) and is
inert in \(K\).  At its degree-two place every rational
quadratic \(\psi_{s,t}\) has Frobenius value one.
The cubic character also has value one: in the
\(S_3\) splitting field, a rational Frobenius outside
the quadratic subgroup is a reflection, whose square
is one.  The original diagonal trace is consequently
the number \(b\) in (G124.8), and \(b>3\).
At actual Type I points the normalized automorphic
trace has complex absolute value at most three.
The degree-zero trace difference therefore proves M5
by exactly the Type I argument of AC-G122.

All five model conditions now hold, so the original
\(\mathrm{LemmaM}\), \(\mathrm{LemmaM'}\) and strict
integral extraction give the first two inequalities
of (G124.9), with \(v_{\mathcal O_0}=e_{\mathcal O_0}v_{269}\).
AC-G111 applies to the same actual CM tower, whose
prime degrees remain \(2\) and \(269\).  The rank-three
determinant is still \(\chi^{2j}\), independent of the
quadratic \(\psi\), so its determinant normalization
and finite-twist removal are unchanged.
The resulting original-\(K\) representation has its
own exact inertia (G124.16) and away type.

AC-G112-local consequently gives principal compact
level \(269N_{s,t}\).  This theorem explicitly permits
\(269\mid N\): at that prime the conductor bound is
\(\max(1,v_{269}(N))\), while \(269N\) is a uniform
upper bound.  It does not require the false literal
old \(\mathrm{LevelAt}\) cyclotomic condition for
the cubic character.  AC-G112-count thus gives one
original-field classifier of size
\(C_1\{N_{s,t}(e+1)\}^{A_1}\), proving (G124.9).
This is where the displayed conductor must change
from \(4\) to the correct row of (G124.3).

Finally the source trace-norm theorem used in
AC-G122 applies to this classifier, with \(Q=121\),
\(z(j)=1\), and the given set of labels.  It gives
one nonzero rational integer with the required
product valuation inequality.  Its logarithmic
bound is at most
\(\#J_0C_1Y^{A_1}(2e\log121+1)\), which is at most
\(5C_1\#J_0Y^{A_1+1}\log11\).
Summing the original-length inequalities and using
\(v_{269}(n)\log269\le\log|n|\) gives exactly
(G124.10).  This completes the proof.
\end{proof}

\paragraph{Scope and geometric dictionary.}
These are four actual original model families with common host
support, under the stated retained native and published inputs.
The result does not assert an original Type II seed in any of
them, transport the AC-G103 class certificate to an arbitrary
prime-to-\(p\) extension, or identify different selected local
components.  The unramified character \(\psi_0\) with
\(e\equiv133\) is not one of the two \(t=1\) rows: there
the compensating \(\tau\) is essential to (G124.13).
No assertion for all characters or all coefficient primes in
H04.4 follows from this finite family.

Geometrically each original model maps, point by point or curve
by curve, to the same normalized ordinary deformation chart on
the host.  The finite-cover and local-system interpretation is
the one already audited for AC-G120--AC-G123 using SGA~1,
Expose~V, sections~4 and~7.  The character identity here is
restriction and tensoring by the stated finite quadratic
characters.  It is not a new motivic realization or a new
cohomological comparison theorem.

\paragraph{Exact proof and source scope.}
The complete AC-G123 composition was independently reread and
reproved before this extension.  The complete native
\texttt{Counts.lean}, \texttt{RhoKApply.lean} and
\texttt{DirichletPrimeToLPartMulLPowerPart.lean} containing
files were freshly read, as were the complete AC-G65 local
component/flag/model text, AC-G109, AC-G117 and AC-G122.
The actual AC-G112-local setup and full local proof, and the
applied counting statement and proof interface, were freshly
checked; its prior full source audit remains retained.
The preceding AC-G121 full native-source audit and the
recorded AC-G123 dependency audits supply the unchanged
\(\mathrm{Resid}\), \(\mathrm{Pair}\), \(\mathrm{PT}\),
strong-seed and uniform-support interfaces.  No new proof
of a deep published input or Lean build is claimed.
